\section{Cycle Extraction} \label{sec:cycle-extraction}

The second category of potential evidence for proictal states was based on cyclical patterns of seizure occurrences and predictions in the time series data.
This section describes the methods used to preprocess the data, handle missing values, extract relevant events, and perform circular statistical analyses to relate cycles in features to the timing of seizures and model predictions.

\subsection{Gap Filling} \label{subsec:gap-filling}
Cycle extraction required gap-free data, so gaps were filled before processing.
To avoid losing data from invalid clips, this was performed on the segment- rather than clip-level, such that non-existing segments were imputed.

The gap-filling procedure preserved the original segment timeline while imputing only short missing intervals ($<\SI{14}{\day}$) in the extracted feature series, as in \textcite{yang_seizure_2024}.
Long gaps ($\geq\SI{14}{\day}$) were treated as structural recording interruptions and were therefore not imputed.
Instead, they were used to split the data into independent continuous chunks, ensuring that feature samples were never drawn across major recording breaks.

Within each short gap, imputation was performed independently per feature column by sampling adjacent segment values in the following manner:
\begin{enumerate}
    \item For a gap of length $L$, a \textbf{pool} of feature value samples was taken from neighboring segments before and after the gap within a window of size $W=\max(L, N_{\min})$, where $N_{\min}=100$ is the minimum sampling span.
    \item Each missing value was \textbf{imputed} with a random value from the pool, sampling with replacement.
    \item Additionally, each sampled value was \textbf{perturbed} by a random multiplicative factor in the interval $[1-\delta,\;1+\delta]$, where $\delta=0.05$ is the maximum distortion percentage. 
\end{enumerate}
This design preserved local feature statistics, avoided unrealistic interpolation over long outages, and introduced slight stochastic variation to reduce deterministic copy artifacts in the imputed sequences.


\subsection{Events} \label{subsec:cycle-events}
Three types of events were evaluated: seizure onsets and \glspl{fp} predicted by both models.
Only \glspl{fp} from valid clips in the test set were included, determined using the optimal threshold (\cref{subsec:threshold-optimization}).
Due to the low number of seizures ($\leq$10) in the test sets of some patients, and assuming stable phase locking of seizures to cycles over time, seizures were drawn from the entire dataset to enable statistical analysis.


\subsection{Outlier Handling} \label{subsec:cycle-outlier}
Feature-based phase locking can be affected by artifacts remaining in the \gls{eeg} signal, including myogenic activity and technical noise.
To attenuate the impact of such artifacts, extreme feature values were clipped at the \SI{99.999}{\percent} quantile.
A range of quantile thresholds was systematically evaluated across all patients and all feature types, and the chosen threshold was identified through visual inspection as effectively suppressing extreme outliers while retaining the vast majority of the data.
Subsequent to clipping, all features were standardized via z-score normalization to standardize their distributions.

\subsection{Filtering} \label{subsec:cycle-filtering}
A non-causal, multidien filter was applied to each continuous feature chunk to isolate slow oscillatory components.
Specifically, a zero-phase 10th order Butterworth bandpass filter in \gls{sos} form was used with a period range of 5--50 days, and a sampling frequency corresponding to the segment duration.
This process extracted multidien dynamics while removing short-term fluctuations.

\subsection{Phase Locking} \label{subsec:phase-locking}
Based on the filtered features, the phases at event times, \glspl{plv}, and mean angles were computed using the Hilbert transform.
The \gls{plv} quantifies the consistency of event timing relative to the phase of an oscillatory process, ranging from 0 (no phase preference) to 1 (perfect phase locking).
This is a standard tool in circular statistics and neuroscience for quantifying phase-event relationships.

\paragraph{Computation of PLV}
\begin{enumerate}
    \item \textbf{Analytic Signal:} Given a real-valued time series $x(t)$ (e.g., a filtered feature), the analytic signal $z(t)$ is computed using the Hilbert transform:
    \begin{equation} \label{eq:analytic-signal}
        z(t) = x(t) + i \cdot \mathcal{H}[x(t)]
    \end{equation}
    where $\mathcal{H}[x(t)]$ is the Hilbert transform of $x(t)$.

    \item \textbf{Event Phase Extraction:}
    The instantaneous phase is then:
    \begin{equation} \label{eq:instantaneous-phase}
        \phi(t) = \arg(z(t)).
    \end{equation}
    For each event (e.g., a seizure / \gls{fp}), extract the instantaneous phase at the event times where $N$ is the number of events: $\{\phi_{e_1}, \phi_{e_2}, \ldots, \phi_{e_N}\}$.

    \item \textbf{Mean Resultant Vector:} The mean resultant vector is
    \begin{equation} \label{eq:mean-resultant-vector}
        R = \frac{1}{N} \sum_{k=1}^N e^{i \phi_{e_k}}.
    \end{equation}

    \item \textbf{PLV Calculation:} The \gls{plv} is the magnitude of the mean resultant vector:
    \begin{equation} \label{eq:plv}
        \mathrm{PLV} = |R| = \left| \frac{1}{N} \sum_{k=1}^N e^{i \phi_{e_k}} \right|
    \end{equation}

    \item \textbf{Mean Angle:} The argument of $R$ gives the preferred phase:
    \begin{equation} \label{eq:plv-mean-angle}
        \mu = \arg(R)
    \end{equation}
\end{enumerate}

\subsection{Rayleigh Test} \label{subsec:rayleigh}
The statistical significance of phase locking was assessed using the Rayleigh test for each feature, which tests the null hypothesis that the event phases are uniformly distributed.

\subsection{Circular Comparison} \label{subsec:circular-comparison}
To determine whether two groups of circular data, i.e., event phases from two event types, had significantly different mean directions, a non-parametric permutation test\footnote{The Watson-Wheeler test, which tests for equality of circular distributions, was also considered as an alternative but found to be overly sensitive to event phase variability.} was used.
This is appropriate for circular data, where standard linear tests are not valid due to the periodic nature of the data.
The null hypothesis that both samples are drawn from populations with the same mean direction was tested.

Let $\{\phi_1^{(1)}, \ldots, \phi_{n_1}^{(1)}\}$ and $\{\phi_1^{(2)}, \ldots, \phi_{n_2}^{(2)}\}$ be the two samples of event phases (in radians), with sample sizes $n_1$ and $n_2$.

\paragraph{Observed Statistic:}
First, compute the circular mean for each group:
\begin{equation} \label{eq:circ-mean}
    \bar{\phi}^{(j)} = \arg\left( \sum_{k=1}^{n_j} e^{i \phi_k^{(j)}} \right), \quad j=1,2
\end{equation}
The observed difference in mean direction is then
\begin{equation} \label{eq:delta-obs}
    \Delta_{\text{obs}} = \mathrm{wrap}(\bar{\phi}^{(1)} - \bar{\phi}^{(2)})
\end{equation}
where $\mathrm{wrap}(\cdot)$ denotes wrapping the angle to the interval $[-\pi, \pi)$.

\paragraph{Variance and Test Statistic:}
To account for the concentration of the data (i.e., how tightly phases are clustered; similar to variance in non-circular statistics), the test uses an estimate of variance based on the mean resultant length $R_j$ for each group:
\begin{equation} \label{eq:mean-resultant-length}
    R_j = \frac{1}{n_j} \left| \sum_{k=1}^{n_j} e^{i \phi_k^{(j)}} \right|
\end{equation}
The variance of the mean direction for each group is approximated as
\begin{equation} \label{eq:var-mean-dir}
    \mathrm{Var}_j = \frac{1 - R_j}{n_j (R_j^2 + \epsilon)}
\end{equation}
where $\epsilon$ is a small constant to avoid division by zero.
The test statistic is then
\begin{equation} \label{eq:circ-tstat}
    T_{\text{obs}} = \frac{|\Delta_{\text{obs}}|}{\sqrt{\mathrm{Var}_1 + \mathrm{Var}_2}}
\end{equation}
This statistic is analogous to a Welch's t-test, but adapted for circular data.

\paragraph{Permutation Procedure:}
To generate the null distribution, the two samples were pooled and randomly reassigned into two groups of sizes $n_1$ and $n_2$, without regard to their original labels.
For each permutation, the difference in circular means and the corresponding test statistic $T_{\text{perm}}$ were computed as above.
This process was repeated a large number of times ($N_{perm}=5000$), yielding a distribution of $T_{\text{perm}}$ under the null hypothesis.

\paragraph{P-value Calculation:}
The two-sided p-value was calculated as the proportion of permutations where the permuted test statistic exceeded or equaled the observed test statistic:
\begin{equation} \label{eq:circ-pval}
    p = \frac{1 + \sum_{i=1}^{N_{\text{perm}}} \mathbb{I}(T_{\text{perm},i} \geq T_{\text{obs}})}{1 + N_{\text{perm}}}
\end{equation}
where $\mathbb{I}(\cdot)$ is the indicator function and $N_{\text{perm}}$ is the number of permutations.

\paragraph{Bootstrap Confidence Intervals:}
In addition to the permutation test, bootstrap resampling was used to estimate confidence intervals for the mean directions and their difference.
This involved repeatedly resampling each group with replacement, recalculating the circular means, and determining the empirical percentiles of the resulting distributions.

\paragraph{Interpretation:}
A small p-value indicates that the observed difference in mean direction is unlikely to have arisen by chance, suggesting a significant difference in the phase preference between the two groups.
The bootstrap confidence intervals provide an estimate of the uncertainty in the mean directions and their difference.

\paragraph{Usage:} The permutation test was applied to all pairs of events: seizures vs.\ \gls{cnn} \glspl{fp}, seizures vs.\ ensemble \glspl{fp}, and \gls{cnn} \glspl{fp} vs.\ ensemble \glspl{fp}.


\subsection{False Discovery Control} \label{subsec:fdr}
Since the Rayleigh and circular tests were performed for each feature (15 p-values per test), results were corrected for \gls{fdr} using the Benjamini-Hochberg procedure, yielding \gls{fdr}-adjusted p-values (also referred to as $q$-values) for each test.
